\begin{textsample}{POS Dim 3 – human – Score 7.00 – mg/mg\_ef\_er\_3\_p2.txt}  \label{ex:f3_pos_005}
Nesse processo formativo, na concepção de Educação do componente Ensino Religioso, o educando é o centro, como sujeito do processo \textbf{educativo}, e deve ser respeitado em sua singularidade, e essa centralidade deve estar expressa no Projeto Pedagógico da escola.
Da mesma forma, o Currículo Referência de Minas Gerais, ao \textbf{definir} a base dos diversos componentes, e aqui do componente Ensino Religioso, propicia que as escolas e suas comunidades \textbf{educativas} se debrucem sobre esse \textbf{documento}, a partir do seu território, avançando na reflexão sobre o planejamento e a execução das aulas : o que ensinar, como ensinar e quando ensinar.
Esse trabalho representa uma \textbf{etapa} imprescindível, pois é preciso adequar o currículo e ressignificá -lo diante das situações regionais e locais e dos sujeitos do processo \textbf{educativo}.
Nessa perspectiva, a \textbf{proposta} \textbf{curricular} do componente Ensino Religioso no CRMG compreende -se como uma “ prática na qual se estabelece um diálogo, por assim dizer, entre agentes sociais, elementos técnicos, alunos que reagem frente a ele, professores que o modelam [... ] ”. ( SACRISTAN, 2000, p. 16 ).

% matched lemmas: curricular, definir, documento, educativo, etapa, proposta
\end{textsample}
